%File: anonymous-submission-latex-2026.tex
\documentclass[letterpaper]{article} % DO NOT CHANGE THIS
\usepackage[submission]{aaai2026}  % DO NOT CHANGE THIS
\usepackage{times}  % DO NOT CHANGE THIS
\usepackage{helvet}  % DO NOT CHANGE THIS
\usepackage{courier}  % DO NOT CHANGE THIS
\usepackage[hyphens]{url}  % DO NOT CHANGE THIS
\usepackage{graphicx} % DO NOT CHANGE THIS
\urlstyle{rm} % DO NOT CHANGE THIS
\def\UrlFont{\rm}  % DO NOT CHANGE THIS
\usepackage{natbib}  % DO NOT CHANGE THIS AND DO NOT ADD ANY OPTIONS TO IT
\usepackage{caption} % DO NOT CHANGE THIS AND DO NOT ADD ANY OPTIONS TO IT
\frenchspacing  % DO NOT CHANGE THIS
\setlength{\pdfpagewidth}{8.5in} % DO NOT CHANGE THIS
\setlength{\pdfpageheight}{11in} % DO NOT CHANGE THIS
%
% These are recommended to typeset algorithms but not required. See the subsubsection on algorithms. Remove them if you don't have algorithms in your paper.
\usepackage{algorithm}

% Add more
\usepackage{subcaption}
\usepackage{algpseudocode} 
\usepackage{amsmath}
\usepackage{multirow}
\usepackage{amsfonts} % For \mathbb
\usepackage{makecell} % Add in preamble

%
% These are recommended to typeset listings but not required. See the subsubsection on listing. Remove this block if you don't have listings in your paper.
\usepackage{newfloat}
\usepackage{listings}
\DeclareCaptionStyle{ruled}{labelfont=normalfont,labelsep=colon,strut=off} % DO NOT CHANGE THIS
\lstset{%
	basicstyle={\footnotesize\ttfamily},% footnotesize acceptable for monospace
	numbers=left,numberstyle=\footnotesize,xleftmargin=2em,% show line numbers, remove this entire line if you don't want the numbers.
	aboveskip=0pt,belowskip=0pt,%
	showstringspaces=false,tabsize=2,breaklines=true}
\floatstyle{ruled}
\newfloat{listing}{tb}{lst}{}
\floatname{listing}{Listing}
%
% Keep the \pdfinfo as shown here. There's no need
% for you to add the /Title and /Author tags.
\pdfinfo{
/TemplateVersion (2026.1)
}

% DISALLOWED PACKAGES
% \usepackage{authblk} -- This package is specifically forbidden
% \usepackage{balance} -- This package is specifically forbidden
% \usepackage{color (if used in text)
% \usepackage{CJK} -- This package is specifically forbidden
% \usepackage{float} -- This package is specifically forbidden
% \usepackage{flushend} -- This package is specifically forbidden
% \usepackage{fontenc} -- This package is specifically forbidden
% \usepackage{fullpage} -- This package is specifically forbidden
% \usepackage{geometry} -- This package is specifically forbidden
% \usepackage{grffile} -- This package is specifically forbidden
% \usepackage{hyperref} -- This package is specifically forbidden
% \usepackage{navigator} -- This package is specifically forbidden
% (or any other package that embeds links such as navigator or hyperref)
% \indentfirst} -- This package is specifically forbidden
% \layout} -- This package is specifically forbidden
% \multicol} -- This package is specifically forbidden
% \nameref} -- This package is specifically forbidden
% \usepackage{savetrees} -- This package is specifically forbidden
% \usepackage{setspace} -- This package is specifically forbidden
% \usepackage{stfloats} -- This package is specifically forbidden
% \usepackage{tabu} -- This package is specifically forbidden
% \usepackage{titlesec} -- This package is specifically forbidden
% \usepackage{tocbibind} -- This package is specifically forbidden
% \usepackage{ulem} -- This package is specifically forbidden
% \usepackage{wrapfig} -- This package is specifically forbidden
% DISALLOWED COMMANDS
% \nocopyright -- Your paper will not be published if you use this command
% \addtolength -- This command may not be used
% \balance -- This command may not be used
% \baselinestretch -- Your paper will not be published if you use this command
% \clearpage -- No page breaks of any kind may be used for the final version of your paper
% \columnsep -- This command may not be used
% \newpage -- No page breaks of any kind may be used for the final version of your paper
% \pagebreak -- No page breaks of any kind may be used for the final version of your paper
% \pagestyle -- This command may not be used
% \tiny -- This is not an acceptable font size.
% \vspace{- -- No negative value may be used in proximity of a caption, figure, table, section, subsection, subsubsection, or reference
% \vskip{- -- No negative value may be used to alter spacing above or below a caption, figure, table, section, subsection, subsubsection, or reference

\setcounter{secnumdepth}{0} %May be changed to 1 or 2 if section numbers are desired.

% The file aaai2026.sty is the style file for AAAI Press
% proceedings, working notes, and technical reports.
%

% Title

% Your title must be in mixed case, not sentence case.
% That means all verbs (including short verbs like be, is, using,and go),
% nouns, adverbs, adjectives should be capitalized, including both words in hyphenated terms, while
% articles, conjunctions, and prepositions are lower case unless they
% directly follow a colon or long dash
\title{Technical Appendix for\\ SHeRL-FL: When Representation Learning Meets Split Learning in Hierarchical Federated Learning}
\author{
    %Authors
    % All authors must be in the same font size and format.
    Written by AAAI Press Staff\textsuperscript{\rm 1}\thanks{With help from the AAAI Publications Committee.}\\
    AAAI Style Contributions by Pater Patel Schneider,
    Sunil Issar,\\
    J. Scott Penberthy,
    George Ferguson,
    Hans Guesgen,
    Francisco Cruz\equalcontrib,
    Marc Pujol-Gonzalez\equalcontrib
}
\affiliations{
    %Afiliations
    \textsuperscript{\rm 1}Association for the Advancement of Artificial Intelligence\\
    % If you have multiple authors and multiple affiliations
    % use superscripts in text and roman font to identify them.
    % For example,

    % Sunil Issar\textsuperscript{\rm 2},
    % J. Scott Penberthy\textsuperscript{\rm 3},
    % George Ferguson\textsuperscript{\rm 4},
    % Hans Guesgen\textsuperscript{\rm 5}
    % Note that the comma should be placed after the superscript

    1101 Pennsylvania Ave, NW Suite 300\\
    Washington, DC 20004 USA\\
    % email address must be in roman text type, not monospace or sans serif
    proceedings-questions@aaai.org
%
% See more examples next
}

%Example, Single Author, ->> remove \iffalse,\fi and place them surrounding AAAI title to use it
\iffalse
\title{My Publication Title --- Single Author}
\author {
    Author Name
}
\affiliations{
    Affiliation\\
    Affiliation Line 2\\
    name@example.com
}
\fi

\iffalse
%Example, Multiple Authors, ->> remove \iffalse,\fi and place them surrounding AAAI title to use it
\title{My Publication Title --- Multiple Authors}
\author {
    % Authors
    First Author Name\textsuperscript{\rm 1},
    Second Author Name\textsuperscript{\rm 2},
    Third Author Name\textsuperscript{\rm 1}
}
\affiliations {
    % Affiliations
    \textsuperscript{\rm 1}Affiliation 1\\
    \textsuperscript{\rm 2}Affiliation 2\\
    firstAuthor@affiliation1.com, secondAuthor@affilation2.com, thirdAuthor@affiliation1.com
}
\fi


% REMOVE THIS: bibentry
% This is only needed to show inline citations in the guidelines document. You should not need it and can safely delete it.
\usepackage{bibentry}
% END REMOVE bibentry

\begin{document}
\maketitle

\section{Detailed Experimental Setup}
    \textbf{Hardware and Software Environment.} All experiments were conducted on a personal workstation equipped with an NVIDIA GeForce RTX 3080 GPU (10\,GB VRAM), AMD Ryzen 9 5900X CPU (12 cores, 3.7\,GHz), and 32\,GB RAM. The operating system was Ubuntu 22.04 LTS, running Python 3.9.12. The implementation was based on PyTorch 2.7.0+cu128 with Torchvision 0.22.0+cu128, Torchaudio 2.7.0+cu128, and CUDA 12.8. Key supporting libraries include NumPy 2.0.2, Scikit-learn 1.6.1, Pandas 2.2.3, Matplotlib 3.9.4, and HuggingFace Transformers 4.53.1. Additional federated learning components were implemented with Flower 1.18.0, Flower Datasets 0.5.0, and PyTorch Lightning 2.5.2. Random seeds were fixed across all runs, and deterministic \texttt{cuDNN} settings were enabled to ensure reproducibility.

    \noindent\textbf{Datasets.} We evaluate our method on four benchmark datasets:
        \begin{itemize}
            \item CIFAR-10 and CIFAR-100 ~\cite{krizhevsky2009learning} consist of 60,000 color images of size $32 \times 32$ in 10 and 100 classes, respectively. Both datasets are split into 50,000 training images and 10,000 test images. Images are normalized using dataset-specific mean and standard deviation, and augmented via random horizontal flip and random crop.
            \item HAM10000~\cite{tschandl2018ham10000} is a dermatoscopic image dataset containing 10,015 images across 7 diagnostic categories. We follow the standard train/validation/test split (70\%/10\%/20\%) with stratification by class. Images are normalized using dataset-specific mean and standard deviation, and augmented via random horizontal flip and random crop.
            \item ISIC-2018~\cite{8363547} contains 10,015 dermatoscopic images for multi-class skin lesion classification. We adopt the official training/validation/test partitions. Preprocessing follows the same protocol as HAM10000.
        \end{itemize}
        For CIFAR-10, CIFAR-100, and HAM10000 classification tasks, we employ cross-entropy loss with class weights to address imbalance (notably in HAM10000). For ISIC-2018 segmentation, we utilize focal loss to handle severe foreground–background imbalance. All datasets are loaded using \texttt{torchvision.datasets} (for CIFAR) or custom \texttt{torch.utils.data.Dataset} loaders with \texttt{PIL} (for HAM10000 and ISIC-2018).
    
    \noindent\textbf{Model Architectures.} We design dataset-specific architectures for both classification and segmentation tasks, selecting backbones appropriate to each benchmark. For CIFAR-10 classification, we employ AlexNet as the backbone; for CIFAR-100, we adopt ResNet-18; for HAM10000 classification, we use ResNet-50 to capture the high variability of dermatoscopic images; and for ISIC-2018 segmentation, we utilize a U-Net with a ResNet-50 encoder. Tables~\ref{tab:model_cifar10}, \ref{tab:model_cifar100}, \ref{tab:model_ham10000}, and \ref{tab:model_isic2018} present the model partition across tiers for CIFAR-10, CIFAR-100, HAM10000, and ISIC-2018, respectively.
        \begin{table}
            \centering
            \small
            \begin{tabular}{l|l|l}
            \hline
            \textbf{Method} & \textbf{Tiers} & \textbf{Architecture} \\
            \hline
            \multirow{3}{*}{Traditional FL \& HierFL} 
            & Clients & AlexNet (full) \\
            & Edge & N/A \\
            & Cloud & N/A \\
            \hline
            \multirow{3}{*}{SplitFL} 
            & Clients & AlexNet Conv1 \\
            & Edge & \makecell[l]{AlexNet Conv2--Conv5 \\ + FC} \\
            & Cloud & N/A \\
            \hline
            \multirow{3}{*}{HSFL} 
            & Clients & AlexNet Conv1 \\
            & Edge & AlexNet Conv2--Conv4 \\
            & Cloud & \makecell[l]{AlexNet Conv5 \\ + FC} \\
            \hline
            \multirow{3}{*}{Ours} 
            & Clients & AlexNet Conv1 \\
            & Edge & AlexNet Conv2--Conv5 \\
            & Cloud & FC \\ \hline
            \end{tabular}
            \caption{Model partition across tiers for each method on CIFAR-10.}
            \label{tab:model_cifar10}
        \end{table}
        
        \begin{table}
            \centering
            \small
            \begin{tabular}{l|l|l}
            \hline
            \textbf{Method} & \textbf{Tiers} & \textbf{Architecture} \\
            \hline
            \multirow{3}{*}{Traditional FL \& HierFL} 
            & Clients & ResNet-18 (full) \\
            & Edge & N/A \\
            & Cloud & N/A \\
            \hline
            \multirow{3}{*}{SplitFL} 
            & Clients & \makecell[l]{ResNet-18 Conv1 \\ + Layer1--Layer2} \\
            & Edge & \makecell[l]{ResNet-18 Layer3--Layer4 \\ + AvgPool + FC} \\
            & Cloud & N/A \\
            \hline
            \multirow{3}{*}{HSFL} 
            & Clients & \makecell[l]{ResNet-18 Conv1 \\ + Layer1} \\
            & Edge & ResNet-18 Layer2--Layer3 \\
            & Cloud & \makecell[l]{ResNet-18 Layer3--Layer4 \\ + FC} \\
            \hline
            \multirow{3}{*}{Ours} 
            & Clients & \makecell[l]{ResNet-18 Conv1 \\ + Layer1} \\
            & Edge & ResNet-18 Layer2--Layer4 \\
            & Cloud & FC \\ \hline
            \end{tabular}
            \caption{Model partition across tiers for each method on CIFAR-100.}
            \label{tab:model_cifar100}
        \end{table}
        
        \begin{table}
            \centering
            \small
            \begin{tabular}{l|l|l}
            \hline
            \textbf{Method} & \textbf{Tiers} & \textbf{Architecture} \\
            \hline
            \multirow{3}{*}{Traditional FL \& HierFL} 
            & Clients & ResNet-50 (full) \\
            & Edge & N/A \\
            & Cloud & N/A \\
            \hline
            \multirow{3}{*}{SplitFL} 
            & Clients & \makecell[l]{ResNet-50 Conv1--Layer1 \\ --Layer2} \\
            & Edge & \makecell[l]{ResNet-50 Layer3--Layer4 \\ + AvgPool + FC} \\
            & Cloud & N/A \\
            \hline
            \multirow{3}{*}{HSFL} 
            & Clients & ResNet-50 Conv1--Layer1 \\
            & Edge & ResNet-50 Layer2--Layer3 \\
            & Cloud & \makecell[l]{ResNet-50 Layer4 \\ + AvgPool + FC} \\
            \hline
            \multirow{3}{*}{Ours} 
            & Clients & ResNet-50 Conv1--Layer1 \\
            & Edge & \makecell[l]{ResNet-50 Layer2--Layer4 \\ + AvgPool} \\
            & Cloud & FC \\ \hline
            \end{tabular}
            \caption{Model partition across tiers for each method on HAM10000.}
            \label{tab:model_ham10000}
        \end{table}
        
        \begin{table}
            \centering
            \small
            \begin{tabular}{l|l|l}
            \hline
            \textbf{Method} & \textbf{Tiers} & \textbf{Architecture} \\
            \hline
            \multirow{3}{*}{Traditional FL \& HierFL} 
            & Clients & U-Net (full) \\
            & Edge & N/A \\
            & Cloud & N/A \\
            \hline
            \multirow{3}{*}{SplitFL} 
            & Clients & U-Net Encoder Conv1 \\
            & Edge & \makecell[l]{U-Net Encoder Layer2--Layer4 \\ + U-Net Decoder} \\
            & Cloud & N/A \\
            \hline
            \multirow{3}{*}{HSFL} 
            & Clients & U-Net Encoder Conv1 \\
            & Edge & U-Net Encoder Layer2--Layer3 \\
            & Cloud & \makecell[l]{U-Net Encoder Layer4 \\ + U-Net Decoder} \\
            \hline
            \multirow{3}{*}{Ours} 
            & Clients & U-Net Encoder Conv1 \\
            & Edge & U-Net Encoder Layer2--Layer4 \\
            & Cloud & U-Net Decoder \\ \hline
            \end{tabular}
            \caption{Model partition across tiers for each method on ISIC-2018.}
            \label{tab:model_isic2018}
        \end{table}




    \noindent\textbf{Training Configuration.} For all datasets, we employed the Adam optimizer with an initial learning rate of $1 \times 10^{-4}$, weight decay of $5 \times 10^{-4}$, and a cosine annealing learning rate scheduler. Batch size was fixed at $32$ for all experiments to ensure consistent memory usage and fair comparison across datasets.

    We used cross-entropy loss for all classification tasks. For segmentation (ISIC-2018), we adopted focal loss with $\gamma=2$ to mitigate class imbalance. All experiments were trained for $200$ communication rounds, with model selection based on the best validation performance.
    
    \noindent\textbf{Federated Learning Configuration.} We consider two federation topologies.
    \begin{itemize}
        \item In the client–cloud configuration (used for Traditional FL and SplitFL), all $N_c = 200$ clients communicate directly with a single cloud server. Aggregation is performed at the cloud level after each round.
        \item In the client–edge–cloud configuration (used for HierFL, HSFL, and Ours), $N_c = 200$ clients are evenly partitioned across $N_e = 10$ edge servers ($20$ clients per edge). Each edge server performs local aggregation of its connected clients before forwarding updates to the cloud. The cloud aggregates the updates received from all edges to update the global model.
    \end{itemize} 
    All configurations are trained for $R = 200$ communication rounds. In each round, only $10\%$ of clients are randomly selected to participate in the training process. No compression, quantization, or sparsification is applied, ensuring a consistent basis for performance and communication analysis.

    \noindent\textbf{Evaluation Protocol.} Model performance is evaluated on a held-out test set for each dataset under both IID and non-IID client data distributions for the classification task and IID for the segmentation task:
        \begin{itemize}
            \item For classification tasks (CIFAR-10, CIFAR-100, HAM10000), the primary evaluation metric is the macro-averaged F1-score to account for class imbalance. 
            \item For segmentation tasks (ISIC-2018), evaluation focuses on pixel-level performance using the Dice Coefficient (DICE) and Intersection-over-Union (IoU).
            \item For computational cost analysis, we record the average RAM usage of each client throughout the training process, expressed in MB. 
            \item Communication cost is measured as the total data volume exchanged between devices over all communication rounds, averaged per round, and reported in GB.
        \end{itemize}
\section{Extra Results}
    \begin{table*}
        \centering
        \begin{tabular}{l|cccccccc|c}
        \hline
        \textbf{Method} &
        \begin{tabular}[c]{@{}c@{}}\textbf{Client}\\ \textbf{smashed}\\ \textbf{upload}\end{tabular} &
        \begin{tabular}[c]{@{}c@{}}\textbf{Edge}\\ \textbf{smashed}\\ \textbf{upload}\end{tabular} &
        \begin{tabular}[c]{@{}c@{}}\textbf{Client}\\ \textbf{model}\\ \textbf{upload}\end{tabular} &
        \begin{tabular}[c]{@{}c@{}}\textbf{Client}\\ \textbf{model}\\ \textbf{download}\end{tabular} &
        \begin{tabular}[c]{@{}c@{}}\textbf{Edge}\\ \textbf{model}\\ \textbf{upload}\end{tabular} &
        \begin{tabular}[c]{@{}c@{}}\textbf{Edge}\\ \textbf{model}\\ \textbf{download}\end{tabular} &
        \begin{tabular}[c]{@{}c@{}}\textbf{Gradient}\\ \textbf{from}\\ \textbf{cloud}\end{tabular} &
        \begin{tabular}[c]{@{}c@{}}\textbf{Gradient}\\ \textbf{from}\\ \textbf{edge}\end{tabular} &
        \textbf{Total (GB)} \\
        \hline
        FedAvg & -- & -- & 351.40 & 3,513.95 & -- & -- & -- & -- & 3,865.35 \\
        FedSGD & -- & -- & 351.40 & 3,513.95 & -- & -- & -- & -- & 3,865.35 \\
        FedNova & -- & -- & 351.40 & 3,513.95 & -- & -- & -- & -- & 3,865.35 \\
        FedProx & -- & -- & 351.40 & 3,513.95 & -- & -- & -- & -- & 3,865.35 \\ \hline
        HierFL & -- & -- & 351.40 & 3,865.34 & 175.70 & -- & -- & -- & 4,392.44 \\ \hline
        SplitFed & 6.87 & -- & 1.09 & 216.82 & -- & -- & 3.73 & -- & 228.51 \\ \hline
        HSFL (A) & 54.93 & 13.74 & 0.22 & 0.93 & 1.34 & 1.34 & 12.22 & 31.25 & 115.97 \\
        HSFL (B) & 54.93 & 13.74 & 0.11 & 0.48 & 0.67 & 0.67 & 12.22 & 31.25 & 114.07 \\
        HSFL (C) & 54.93 & 13.74 & 0.06 & 0.24 & 0.34 & 0.34 & 12.22 & 31.25 & 113.12 \\
        HSFL (D) & 54.93 & 13.74 & 0.03 & 0.12 & 0.17 & 0.17 & 12.22 & 31.25 & 112.62 \\
        HSFL (E) & 54.93 & 13.74 & 0.015 & 0.06 & 0.085 & 0.085 & 12.22 & 31.25 & 112.35 \\ \hline
        \textbf{Ours (A)} & \textbf{54.93} & \textbf{13.74} & \textbf{0.22} & \textbf{0.93} & \textbf{1.34} & \textbf{1.34} & -- & \textbf{31.25} & \textbf{103.75} \\
        \textbf{Ours (B)} & \textbf{54.93} & \textbf{13.74} & \textbf{0.11} & \textbf{0.48} & \textbf{0.67} & \textbf{0.67} & -- & \textbf{31.25} & \textbf{101.85} \\
        \textbf{Ours (C)} & \textbf{54.93} & \textbf{13.74} & \textbf{0.06} & \textbf{0.24} & \textbf{0.34} & \textbf{0.34} & -- & \textbf{31.25} & \textbf{100.90} \\
        \textbf{Ours (D)} & \textbf{54.93} & \textbf{13.74} & \textbf{0.03} & \textbf{0.12} & \textbf{0.17} & \textbf{0.17} & -- & \textbf{31.25} & \textbf{100.40} \\
        \textbf{Ours (E)} & \textbf{54.93} & \textbf{13.74} & \textbf{0.015} & \textbf{0.06} & \textbf{0.085} & \textbf{0.085} & -- & \textbf{31.25} & \textbf{100.15} \\
        \hline
        \end{tabular}
        \caption{Communication overhead (GB) on CIFAR-100 with ResNet-18 over 200 rounds. 
        (A) $(t_1=5, t_2=10)$, (B) $(t_1=10, t_2=20)$, (C) $(t_1=25, t_2=50)$, (D) $(t_1=50, t_2=100)$, and (E) $(t_1=100, t_2=200)$.}
        \label{tab:comm_overhead_extended}
    \end{table*}
    In extended configurations \textbf{(D)} $(t_1=50, t_2=100)$ and \textbf{(E)} $(t_1=100, t_2=200)$, which simulate highly infrequent aggregation cycles, our method continues to demonstrate substantial communication savings over HSFL while maintaining competitive performance. As shown in Table~\ref{tab:comm_overhead_extended}, the total communication drops to $\sim100.15$~GB in \textbf{(E)} for our method versus $\sim112.35$~GB for HSFL, achieving $\sim10.8\%$ additional reduction compared to the baseline. Although convergence latency increases due to less frequent synchronization, the accuracy degradation from \textbf{C} $\rightarrow$ \textbf{D} $\rightarrow$ \textbf{E} remains modest, suggesting that the edge projection head and enhanced feature alignment effectively mitigate model drift. These results indicate that configurations \textbf{D/E} are well-suited for bandwidth-constrained deployments where communication efficiency is prioritized over rapid convergence.
    \begin{table}
        \centering
        \begin{tabular}{l|c|c}
        \hline
        \textbf{Methods} & \textbf{IoU (\%)} & \textbf{DICE (\%)} \\ \hline
        FedAvg    & 13.73  & 23.38 \\ 
        FedSGD    & 10.23  & 20.54 \\ 
        FedNova   & 19.77  & 25.62 \\ 
        FedProx   & 22.01  & 29.93 \\ \hline
        SplitFed  & 15.26  & 24.46 \\ \hline
        HierFL    & 15.93  & 25.06 \\ \hline
        HSFL (A)  & 24.45 & 32.20 \\
        HSFL (B)  & 24.12 & 32.05 \\
        HSFL (C)  & 23.88 & 31.74 \\
        HSFL (D)  & 23.55 & 31.50 \\
        HSFL (E)  & 23.20 & 31.25 \\ \hline
        Ours (A)  & 28.30 & 35.10 \\
        Ours (B)  & 27.05 & 34.85 \\
        Ours (C)  & 26.72 & 34.42 \\
        Ours (D)  & 26.40 & 34.15 \\
        Ours (E)  & 26.10 & 33.90 \\ \hline
        \end{tabular}
        \caption{Segmentation performance comparison of different FL methods in terms of IoU and DICE scores. (A) $(t_1=5, t_2=10)$, (B) $(t_1=10, t_2=20)$, (C) $(t_1=25, t_2=50)$, (D) $(t_1=50, t_2=100)$, and (E) $(t_1=100, t_2=200)$.}
        \label{tab:segmentation_extended}
    \end{table}
    In extended segmentation configurations \textbf{(D)} $(t_1=50, t_2=100)$ and \textbf{(E)} $(t_1=100, t_2=200)$, where aggregation becomes significantly less frequent, our method maintains its performance advantage over HSFL despite the reduced synchronization. As presented in Table~\ref{tab:segmentation_extended}, the IoU for our method only drops from $26.72\%$ in \textbf{(C)} to $26.10\%$ in \textbf{(E)}, and DICE from $34.42\%$ to $33.90\%$. In contrast, HSFL shows a larger relative decline, with IoU decreasing from $23.88\%$ to $23.20\%$ and DICE from $31.74\%$ to $31.25\%$. This demonstrates the robustness of our approach to prolonged local training intervals, with the projection head and hierarchical feature alignment effectively mitigating model drift. These results suggest that configurations \textbf{D/E} are practical for segmentation tasks in bandwidth-constrained deployments, offering strong accuracy retention alongside substantial communication savings.

\bibliography{aaai2026}

\end{document}
